\documentclass[11pt,a4paper]{article}
\usepackage[margin=1in]{geometry}
\usepackage{amsmath,amssymb,amsthm}
\usepackage{algorithm}
\usepackage{algpseudocode}
\usepackage{booktabs}
\usepackage{hyperref}

\theoremstyle{definition}
\newtheorem{definition}{\textbf{Definition}}[section]
\newtheorem{theorem}{\textbf{Theorem}}[section]
\newtheorem{lemma}[theorem]{\textbf{Lemma}}
\newtheorem{remark}{\textbf{Remark}}[section]

\title{\textbf{Rigorous Mathematical Specification and Cumulative Mass Filtering of Nucleus (Top-$p$) and Min-$p$ Sampling}}
\author{\textbf{Team Scaibu} \\ \small Email: \texttt{jaydeep.9664920749@gmail.com} \\ \small \textbf{Architecture and Core Engine Team}}
\date{\textbf{October 2026}}

\begin{document}
\maketitle

\begin{abstract}
Fixed $k$ in Top-$k$ sampling is inherently brittle: when the model's predictive distribution is flat, $k$ is artificially restrictive; conversely, when the distribution is highly peaked, $k$ permits tail tokens with negligible probability to leak into the candidate pool. Nucleus (Top-$p$) sampling dynamically adapts candidate pool size by selecting the smallest subset of tokens whose cumulative probability mass exceeds $p \in (0, 1]$. Min-$p$ sampling introduces an adaptive relative probability floor, discarding any token whose probability drops below a threshold fraction of the top token's probability. This document formalizes both truncation operators, derives normalization invariants, details computational complexities, proves probability concentration bounds, and provides a verified step-by-step numerical calculation.
\end{abstract}

\section{Notation and Mathematical Preliminaries}

Let $\mathcal{V} = \{1, 2, \dots, V\}$ denote the vocabulary with cardinality $|\mathcal{V}| = V$. Given unnormalized logits $\mathbf{z} \in \mathbb{R}^V$, the base probability distribution is given by the stable softmax operator:
\begin{equation}
P(v) = \frac{\exp(z_v - \max_{j} z_j)}{\sum_{l=1}^V \exp(z_l - \max_{j} z_j)}, \quad \forall v \in \mathcal{V}
\end{equation}

\begin{table}[h]
\centering
\caption{Mathematical Notation for Nucleus and Min-$p$ Sampling}
\begin{tabular}{lll}
\toprule
\textbf{Symbol} & \textbf{Domain} & \textbf{Semantic Description} \\
\midrule
$V$ & $\mathbb{N}_{\ge 1}$ & Vocabulary size \\
$\mathbf{z}$ & $\mathbb{R}^V$ & Unnormalized token logits \\
$P$ & $\Delta^{V-1}$ & Base softmax probability distribution \\
$(1), (2), \dots, (V)$ & Permutation & Indices sorted descending such that $P((1)) \ge P((2)) \ge \dots \ge P((V))$ \\
$p$ & $(0, 1]$ & Top-$p$ cumulative probability mass threshold \\
$p_{\min}$ & $[0, 1]$ & Min-$p$ scaling threshold relative to mode probability \\
$k^*$ & $\mathbb{N}_{\ge 1}$ & Minimal nucleus cutoff index \\
$\mathcal{V}^{(p)}$ & Subset of $\mathcal{V}$ & Nucleus candidate set of retained tokens \\
$P^{(p)}(v)$ & $\Delta^{|\mathcal{V}^{(p)}|-1}$ & Renormalized probability distribution over nucleus \\
\bottomrule
\end{tabular}
\end{table}

\section{Problem Definition and Core Concepts}

\subsection{Intuitive 3-Level Explanation}
\begin{enumerate}
    \item \textbf{Intuitive (Mental Metaphor):} Imagine a chef picking ingredients for a recipe. If the choice is obvious (like tomato sauce for pizza), you only want the #1 best option. If there are 10 great spices that all fit equally well, you want to consider all 10. Top-$k$ forces you to pick exactly $k$ ingredients every time. Nucleus sampling adjusts: when one option is 99\% sure, it considers only that one; when 50 options are equally plausible, it considers all 50.
    \item \textbf{Engineering Level:} Sort tokens descending by probability. Accumulate probability mass until the sum reaches or exceeds $p$ (typically $p=0.90$). Truncate all remaining tokens and set their probabilities to zero. Additionally, if Min-$p$ is active, drop any token whose probability is less than $p_{\min} \cdot P((1))$, preventing tail hallucinations even when $p$ is large. Renormalize the remaining token probabilities so they sum to 1.
    \item \textbf{Mathematical Level:} Given sorted permutation $\pi = ((1), \dots, (V))$, the nucleus cutoff $k^*$ is:
    \begin{equation}
    k^* = \min \left\{ k \in \{1, \dots, V\} \;\Bigg|\; \sum_{i=1}^k P((i)) \ge p \right\}
    \end{equation}
    Under joint Nucleus and Min-$p$ filtering, token $v \in \mathcal{V}$ is retained if and only if:
    \begin{equation}
    v \in \{ (1), \dots, (k^*) \} \quad \text{and} \quad P(v) \ge p_{\min} \cdot P((1))
    \end{equation}
\end{enumerate}

\section{Algorithm Specification}

\begin{algorithm}[H]
\caption{Nucleus (Top-$p$) and Min-$p$ Sampling Pipeline}
\label{alg:nucleus_sampling}
\begin{algorithmic}[1]
\Require Unnormalized logits $\mathbf{z} \in \mathbb{R}^V$, thresholds $p \in (0, 1]$ and $p_{\min} \in [0, 1]$
\Ensure Sampled token $v^* \in \mathcal{V}$, retained candidate indices $\mathcal{R}$, normalized probabilities $\mathbf{p}_{\text{norm}}$
\State $z_{\max} \gets \max_{j} z_j$
\State $e_j \gets \exp(z_j - z_{\max})$ for all $j \in \{1, \dots, V\}$
\State $S \gets \sum_{j=1}^V e_j$
\State $P(j) \gets e_j / S$ for all $j \in \{1, \dots, V\}$
\State $\pi \gets \text{ArgsortDescending}(P)$ \Comment{$\pi[0]$ has highest probability}
\State $P_{\text{top}} \gets P(\pi[0])$
\State $\tau_{\min} \gets p_{\min} \cdot P_{\text{top}}$
\State $\mathcal{R} \gets \emptyset, \quad \text{cum\_mass} \gets 0.0$
\For{$i = 0$ \textbf{to} $V - 1$}
    \State $v \gets \pi[i]$
    \If{$P(v) < \tau_{\min} \textbf{ and } |\mathcal{R}| > 0$}
        \State \textbf{break}
    \EndIf
    \State $\mathcal{R} \gets \mathcal{R} \cup \{v\}$
    \State $\text{cum\_mass} \gets \text{cum\_mass} + P(v)$
    \If{$\text{cum\_mass} \ge p$}
        \State \textbf{break}
    \EndIf
\EndFor
\State $S_{\text{nucleus}} \gets \sum_{v \in \mathcal{R}} P(v)$
\State $p_{\text{norm}}(v) \gets P(v) / S_{\text{nucleus}}$ for all $v \in \mathcal{R}$
\State Sample $v^* \sim p_{\text{norm}}$
\State \Return $\{v^*, \mathcal{R}, \mathbf{p}_{\text{norm}}\}$
\end{algorithmic}
\end{algorithm}

\section{Theoretical Guarantees: Tail Bound Concentration}

\begin{theorem}[Mass Concentration Guarantee]
Let $P \in \Delta^{V-1}$ be a discrete distribution and let $k^*$ be the nucleus threshold index. The retained probability mass before renormalization satisfies:
\begin{equation}
p \le \sum_{v \in \mathcal{V}^{(p)}} P(v) \le p + \max_{v \in \mathcal{V}} P(v)
\end{equation}
Furthermore, the maximum probability discarded is strictly upper bounded:
\begin{equation}
\max_{v \notin \mathcal{V}^{(p)}} P(v) \le \frac{1 - p}{V - k^*}
\end{equation}
\end{theorem}

\begin{proof}
By definition of the minimal index $k^*$, $\sum_{i=1}^{k^*-1} P((i)) < p$. Adding the single boundary token $(k^*)$ gives $\sum_{i=1}^{k^*} P((i)) = \sum_{i=1}^{k^*-1} P((i)) + P((k^*)) < p + P((k^*)) \le p + \max_v P(v)$.
For discarded tokens $i > k^*$, by sorted order $P((i)) \le P((k^*+1))$. Since the remaining $V - k^*$ tokens sum to at most $1 - p$, their maximum cannot exceed the average remaining mass $\frac{1 - p}{V - k^*}$.
\end{proof}

\section{Complexity Analysis}

\begin{itemize}
    \item \textbf{Time Complexity:} $\mathcal{O}(V + k \log V)$ using a partial selection heap, or $\mathcal{O}(V \log V)$ via full quicksort. For top-$p$ cutoff typically $k^* \ll V$, requiring $\mathcal{O}(V)$ expected runtime.
    \item \textbf{Space Complexity:} $\mathcal{O}(V)$ auxiliary buffer for probabilities and index permutation.
\end{itemize}

\section{Worked Numerical Example}

Consider vocabulary $V = 4$ with logits $\mathbf{z} = [4.0, 3.0, 1.0, 0.0]$, top-$p = 0.85$, and min-$p = 0.20$.
\begin{enumerate}
    \item Subtract max logit $z_{\max} = 4.0$: $\tilde{\mathbf{z}} = [0.0, -1.0, -3.0, -4.0]$.
    \item Exponentiated values: $\mathbf{e} = [1.0, 0.367879, 0.049787, 0.018316]$. Sum $S = 1.435982$.
    \item Softmax probabilities: $P = [0.696388, 0.256187, 0.034671, 0.012755]$.
    \item Sorted permutation: tokens $(0), (1), (2), (3)$.
    \item Top token $(0)$: mass $= 0.696388 < 0.85$. Retain token $0$.
    \item Second token $(1)$: cumulative mass $= 0.696388 + 0.256187 = 0.952575 \ge 0.85$. Retain token $1$ and terminate nucleus search!
    \item Min-$p$ threshold: $\tau_{\min} = 0.20 \times 0.696388 = 0.139278$. Since $P((1)) = 0.256187 > \tau_{\min}$, token $1$ is kept.
    \item Renormalization over $\mathcal{R} = \{0, 1\}$: sum $= 0.952575$.
    \begin{equation}
    P_{\text{norm}}(0) = \frac{0.696388}{0.952575} \approx 0.73106, \quad P_{\text{norm}}(1) = \frac{0.256187}{0.952575} \approx 0.26894
    \end{equation}
\end{enumerate}

\begin{thebibliography}{9}
\bibitem{holtzman2019} A.~Holtzman et al., ``The curious case of neural text degeneration,'' in \emph{International Conference on Learning Representations (ICLR)}, 2020.
\bibitem{minp2024} B.~Nguyen et al., ``Min-p sampling for balanced language model decoding,'' in \emph{arXiv preprint arXiv:2407.01082}, 2024.
\end{thebibliography}

\end{document}
